% SUMMARY AND AUDIT TRAIL
% Target: standalone section for problems/3.2/proof.tex; no new packages.
% Ledger map: setup/restart/continuants §§88--99; energy and spectral program
% §§103--104; sigma ladder and the three targets §§107--110; small-gap law
% §§108--109.  Q6500 supplies the u != -1, restricted-window, and index-shift
% corrections; Q6496 supplies renewal/coprimality; Q6506 supplies the block
% argument and GPRV formulation.  Numerical claims point only to
% DOSSIER_E1_EMPIRICAL.md rather than reproducing its tables.
% Proved here, with complete proofs: companion Casoratian [yes]; exact transfer
% product and restart dictionary [yes]; restricted restart inequality [yes];
% orbit reflection [yes, using the already proved strong-reflection theorem];
% continuant reflection, center vanishing, and root pairing [yes]; renewal
% identity [yes]; adjacent coprimality [yes]; ENERGY-3/2 [yes]; fixed-H
% Chebotarev/Burnside consequence [yes, conditional only on the standard
% Chebotarev theorem]; sigma-ladder implication [yes]; h=2 splitting law [yes].
% Open, and explicitly labelled conjectural/programmatic: E1, ATR, GPRV,
% and GAP-2D-SQRT.  No empirical assertion is used in a proof.
% Constant re-derived: YES.  The ordered in-block pairs introduce a factor 2
% omitted in ledger §103.  The resulting leading constant is 2*sqrt(3), about
% 3.464, NOT 2.45; the ceiling in the number of blocks is included below.
% Discrepancies resolved: (i) §103a supersedes §103 and the dossier's 2.45;
% (ii) §110 supersedes the spec's word "equivalent": ATR and GPRV are stronger
% sufficient targets, while GAP-2D-SQRT is a parallel partial target, and the
% three are not formally equivalent; (iii) this section uses the requested
% index set {1,...,p-2}, hence E_p=F_p-(p-2), rather than the dossier's
% length-(p-1) convention; (iv) N_h(r)=U_{h-1}(r+1), not a literal equality
% with unchanged indices.

\section{The projective orbit and its energy: an unconditional
$p^{3/2}$ theorem and the E1 program}
\label{sec:orbit-energy}

Throughout this section $p\ge7$ is prime and
\[
 P(u)=(2u+1)(17u^2+17u+5)
      =34u^3+51u^2+27u+5.
\]
The Ap\'ery recurrence is
\begin{equation}\label{eq:oe-recurrence}
 (n+1)^3y_{n+1}=P(n)y_n-n^3y_{n-1}.
\end{equation}
We use only the regular indices $0\le n\le p-1$; every nonzero integer
denominator that occurs there is reduced modulo~$p$.

% Ledger sources: §§92, 95, 103; Q6506 §1.
\subsection{The companion solution and the two energy conventions}

Let $b$ be the Ap\'ery solution, $b_0=1$, $b_1=5$, and let $c$ be the
companion solution of~\eqref{eq:oe-recurrence} with $c_0=0$, $c_1=1$.

\begin{lemma}[companion Casoratian]\label{lem:oe-casoratian}
For every $n\ge1$,
\[
 W_n:=b_nc_{n-1}-b_{n-1}c_n=-\frac1{n^3}.
\]
In particular, the two solution vectors are linearly independent modulo~$p$
at every index used below.
\end{lemma}

\begin{proof}
The initial value is $W_1=-1$.  Applying~\eqref{eq:oe-recurrence} to both
solutions and subtracting gives
\[
 b_{n+1}c_n-b_nc_{n+1}
 =\frac{n^3}{(n+1)^3}(b_nc_{n-1}-b_{n-1}c_n)
 =\frac{n^3}{(n+1)^3}W_n.
\]
The left side is $W_{n+1}$, so induction proves the formula.
\end{proof}

Put $I_p=\{1,\ldots,p-2\}$, $M=|I_p|=p-2$, and define the projective
orbit
\[
 \pi(n)=[b_n:c_n]\in\mathbb P^1(\mathbb F_p),\qquad n\in I_p.
\]
For $v\in\mathbb P^1(\mathbb F_p)$ let
\[
 m_p(v)=\#\{n\in I_p:\pi(n)=v\}.
\]
We distinguish the full and off-diagonal energies
\begin{equation}\label{eq:oe-energy-conventions}
 \mathcal F_p=\sum_v m_p(v)^2,
 \qquad
 \mathcal E_p=\sum_v m_p(v)(m_p(v)-1).
\end{equation}
Since $\sum_vm_p(v)=M$, the exact reconciliation is
\begin{equation}\label{eq:oe-energy-reconcile}
 \boxed{\mathcal E_p=\mathcal F_p-(p-2).}
\end{equation}
Also $Z_p:=\{n\in I_p:b_n=0\}$ is the fibre over $[0:1]$, so
$|Z_p|=m_p([0:1])$.

% Ledger sources: §§88--89, 98--99; Q6496 §§1--2; Q6500 §§1--6.
\subsection{Restart continuants and the gap-polynomial dictionary}

For $u\ne-1$ set
\[
 M_u=\begin{pmatrix}0&(u+1)^6\\-1&P(u)\end{pmatrix},
 \qquad
 M_u[X:Z]=[(u+1)^6Z:-X+P(u)Z].
\]
Its determinant is $(u+1)^6$, so $M_u$ is a projective transformation
exactly on this domain, and there $M_u(\infty)=0$.  Define
\[
 x_n=[n^3b_{n-1}:b_n]\in\mathbb P^1(\mathbb F_p).
\]
The recurrence gives $x_{n+1}=M_nx_n$.  Thus $b_z=0$ implies
$x_z=\infty$ (the numerator is nonzero by Lemma~\ref{lem:no-consec}) and,
whenever the next step is in the regular window, $x_{z+1}=0$.  This is the
precise restart mechanism used below.  Define
\begin{equation}\label{eq:oe-U-recurrence}
 U_{-1}(s)=0,\quad U_0(s)=1,\quad U_1(s)=P(s),\qquad
 U_{m+1}(s)=P(s+m)U_m(s)-(s+m)^6U_{m-1}(s).
\end{equation}

\begin{lemma}[exact transfer product]\label{lem:oe-transfer-product}
For $m\ge1$, with
$G_m(s)=M_{s+m-1}\cdots M_s$, one has
\begin{equation}\label{eq:oe-transfer-product}
G_m(s)=
\begin{pmatrix}
 -(s+m)^6U_{m-2}(s+1)&(s+m)^6U_{m-1}(s)\\
 -U_{m-1}(s+1)&U_m(s)
\end{pmatrix}.
\end{equation}
Moreover $\deg U_m=3m$ over~$\mathbb Q$.
\end{lemma}

\begin{proof}
For $m=1$ the formula is the definition of $M_s$.  Left multiplication by
$M_{s+m}$ makes the new top row $(s+m+1)^6$ times the old bottom row, while
the new bottom row is minus the old top row plus $P(s+m)$ times the old
bottom row.  The two entries are exactly those prescribed by
\eqref{eq:oe-U-recurrence}, proving the matrix formula by induction.

If $\lambda_m$ is the leading coefficient of $U_m$, then
$\lambda_0=1$, $\lambda_1=34$, and
$\lambda_{m+1}=34\lambda_m-\lambda_{m-1}$.  The characteristic roots
$17\pm12\sqrt2$ are positive, whence $\lambda_m>0$ and $\deg U_m=3m$.
\end{proof}

Let $N_h(r)$ retain the established positional-gap convention of
Lemma~\ref{lem:gap-poly}: $N_1=1$, $N_2(r)=P(r+1)$, and
\[
 N_{h+1}(r)=P(r+h)N_h(r)-(r+h)^6N_{h-1}(r).
\]

\begin{proposition}[restart and the $N$--$U$ dictionary]
\label{prop:oe-restart-dictionary}
The following statements hold.
\begin{itemize}
\item If $\prod_{j=1}^m(s+j)\ne0$ in $\mathbb F_p$, then the composite of
the $m$ allowed maps starting from $0=[0:1]$ ends at $\infty$ if and only if
$U_m(s)=0$.
\item For every $h\ge1$,
\begin{equation}\label{eq:oe-dictionary}
 \boxed{N_h(r)=U_{h-1}(r+1).}
\end{equation}
On a nonwrapping window, $\pi(r)=\pi(r+h)$ if and only if $N_h(r)=0$.
\item If $1\le z_1<\cdots<z_k\le p-2$ are the zeros of $b$ and
$1\le H\le p-2$, then, for
\begin{equation}\label{eq:oe-Rrel}
 R_{\rm rel}(p,H)=\sum_{m=1}^H
 \#\{1\le s\le p-m-1:U_m(s)=0\},
\end{equation}
one has the restricted restart bound
\begin{equation}\label{eq:oe-restart-bound}
 |Z_p|\le1+R_{\rm rel}(p,H)
 +\left\lfloor\frac{p-3}{H+1}\right\rfloor.
\end{equation}
\end{itemize}
\end{proposition}

\begin{proof}
By Lemma~\ref{lem:oe-transfer-product},
\[
 G_m(s)[0:1]=[(s+m)^6U_{m-1}(s):U_m(s)].
\]
On the displayed open set the product is invertible, so its two coordinates
cannot vanish together.  This proves the first assertion and also explains
why roots at windows containing the forbidden parameter $u=-1$ must be
excluded from the dynamical count.

Both sides of~\eqref{eq:oe-dictionary} have initial values $1,P(r+1)$ and
the same recurrence, so they agree.  For the collision statement, the
solution $y_n=b_rc_n-c_rb_n$ satisfies $y_r=0$ and
$y_{r+1}=b_rc_{r+1}-c_rb_{r+1}=1/(r+1)^3\ne0$ by
Lemma~\ref{lem:oe-casoratian}.  Iterating the recurrence from this zero shows
that $y_{r+h}=0$ exactly when $N_h(r)=0$.  Since
$y_{r+h}=0$ is precisely $\det((b_r,c_r),(b_{r+h},c_{r+h}))=0$, this is the
projective collision condition.

For consecutive zeros put $s_i=z_i+1$ and
$m_i=z_{i+1}-z_i-1$.  Lemma~\ref{lem:no-consec} gives $m_i\ge1$, and the
first assertion gives $U_{m_i}(s_i)=0$.  The pair $(s_i,m_i)$ is unique and
its window is nonwrapping; enlarging the actual range to that in
\eqref{eq:oe-Rrel} is harmless.  Finally
\[
 \sum_{i=1}^{k-1}m_i=z_k-z_1-(k-1)\le p-3.
\]
Thus at most $\lfloor(p-3)/(H+1)\rfloor$ of the $m_i$ exceed $H$; the
additive $1$ pays for the first zero.  This proves~\eqref{eq:oe-restart-bound}.
\end{proof}

\begin{remark}[the one-step shift]\label{rem:oe-one-step-shift}
Immediately after a zero at $z$, the Riccati state has restarted at $0$ at
$s=z+1$.  Hence a positional zero gap $h=z'-z$ contains only $m=h-1$
post-restart maps.  Formula~\eqref{eq:oe-dictionary} records both shifts:
$N_h(z)=U_{h-1}(z+1)$.  Writing ``$N_m=U_m$'' without them changes the
degree and the parity of the center root.
\end{remark}

% Ledger sources: §§89, 94, 103; proof.tex §16 strong reflection.
\subsection{Reflection, center vanishing, and root pairing}

\begin{theorem}[pointwise orbit reflection]\label{thm:oe-orbit-reflection}
For every $n\in I_p$,
\[
 \boxed{\pi(p-1-n)=\pi(n).}
\]
\end{theorem}

\begin{proof}
The strong-reflection theorem, Theorem~\ref{thm:strong-reflection}, applies
to every solution germ of the Ap\'ery recurrence.  Applying it separately
to $b$ and to $c$ gives equality of both homogeneous coordinates, and hence
the projective identity.
\end{proof}

\begin{lemma}[continuant reflection]\label{lem:oe-U-reflection}
For every $m\ge0$,
\begin{equation}\label{eq:oe-U-reflection}
 U_m(-s-m)=(-1)^mU_m(s)
 \qquad\text{in }\mathbb Z[s].
\end{equation}
\end{lemma}

\begin{proof}
$U_m(s)$ is the determinant of the length-$m$ tridiagonal continuant with
\[
 \text{diagonal }P(s),\ldots,P(s+m-1),\qquad
 \text{edge products }(s+1)^6,\ldots,(s+m-1)^6.
\]
Under $s\mapsto-s-m$, reversing the path reverses the edge products and
sends each diagonal entry to its negative, because $P(-1-u)=-P(u)$.
Every determinant term uses
$m-2j$ diagonal entries for some $j$, so it acquires the common sign
$(-1)^{m-2j}=(-1)^m$.  This proves~\eqref{eq:oe-U-reflection}.
\end{proof}

\begin{corollary}[center vanishing]\label{cor:oe-center-vanishing}
If $m$ is odd, then
\[
 U_m(-m/2)=0
 \qquad\text{over }\mathbb Q.
\]
\end{corollary}

\begin{proof}
The point $s=-m/2$ is fixed by $s\mapsto-s-m$.  Lemma
\ref{lem:oe-U-reflection} gives $U_m(-m/2)=-U_m(-m/2)$, and characteristic
zero gives the conclusion.
\end{proof}

\begin{corollary}[root pairing]\label{cor:oe-root-pairing}
Over every field of characteristic different from $2$ in which $U_m$ is
nonzero, the multiset of roots of $U_m$ is stable under $s\mapsto-s-m$.
For odd $m$ the fixed root $-m/2$ is the forced center; every other root
occurs in a two-element orbit.
\end{corollary}

\begin{proof}
Equation~\eqref{eq:oe-U-reflection} preserves vanishing and multiplicity.
The involution has the unique fixed point $-m/2$.
\end{proof}

% Ledger source: §99; Q6496 §2.  The identity was additionally checked there
% by exact symbolic computation for 0 <= m <= 4 and 1 <= g <= 4.
\subsection{Renewal and adjacent coprimality}

Set $G_0(s)$ equal to the identity matrix.

\begin{theorem}[renewal identity]\label{thm:oe-renewal}
For all $m,g\ge0$,
\begin{align}
U_{m+g+1}(s)
={}&U_{m+1}(s)U_g(s+m+1)\notag\\
&-(s+m+1)^6U_m(s)U_{g-1}(s+m+2).
\label{eq:oe-renewal}
\end{align}
\end{theorem}

\begin{proof}
Split the matrix word at the restart edge:
$G_{m+g+1}(s)=G_g(s+m+1)G_{m+1}(s)$.  The lower row of the first factor is
$(-U_{g-1}(s+m+2),U_g(s+m+1))$, while the right column of the second is
$((s+m+1)^6U_m(s),U_{m+1}(s))^{\mathsf T}$ by
Lemma~\ref{lem:oe-transfer-product}.  Their scalar product is the
lower-right entry $U_{m+g+1}(s)$ and is exactly~\eqref{eq:oe-renewal}.
\end{proof}

\begin{theorem}[adjacent coprimality]\label{thm:oe-adjacent-coprime}
For every $m\ge0$,
\[
 \gcd_{\mathbb Q[s]}(U_m,U_{m+1})=1.
\]
\end{theorem}

\begin{proof}
The case $m=0$ is immediate.  Suppose the assertion holds one step earlier.
A common irreducible divisor of $U_m$ and $U_{m+1}$ divides, by
\eqref{eq:oe-U-recurrence}, $(s+m)^6U_{m-1}$.  It cannot divide $U_{m-1}$
by induction, so it must be $s+m$.  But the Ap\'ery recurrence and the same
initial values give $U_m(0)=(m!)^3b_m$; reflection then gives
\[
 U_m(-m)=(-1)^m(m!)^3b_m\ne0
\]
over $\mathbb Q$.  Thus $s+m$ is not a divisor, completing the induction.
\end{proof}

% Ledger sources: §103 and its correction §103a; Q6506 §3.  The collision
% nonvanishing inputs are proof.tex Lemmas gap-poly/nonvanish.  The stronger
% bordered-certificate statement is nv_theorem.tex, Theorem thm:nv-range;
% it is not needed to hide or replace the simpler pair-gap input used here.
\subsection{The unconditional orbit-energy theorem}

For $1\le h\le M-1$, write
\[
 C_h=\#\{1\le r\le M-h:\pi(r)=\pi(r+h)\}.
\]
The Casoratian gives $C_1=0$.  For $h\ge2$, Proposition
\ref{prop:oe-restart-dictionary}, Lemma~\ref{lem:gap-poly}, and
Lemma~\ref{lem:nonvanish} give
\begin{equation}\label{eq:oe-gap-degree-bound}
 C_h\le\deg N_h=3(h-1).
\end{equation}
Theorem~\ref{thm:nv-range} supplies the sharper bordered-certificate
nondegeneracy used elsewhere in the paper; the present argument needs only
the pair-gap nonvanishing in~\eqref{eq:oe-gap-degree-bound}.

\begin{theorem}[orbit energy, ENERGY--$3/2$]
\label{thm:oe-energy-three-halves}
For every prime $p\ge7$ and every integer $2\le H\le M=p-2$,
\begin{equation}\label{eq:oe-block-bound}
 \mathcal F_p\le
 \left\lceil\frac MH\right\rceil
 \bigl(M+3(H-1)(H-2)\bigr).
\end{equation}
Consequently
\begin{equation}\label{eq:oe-three-halves}
 \boxed{\mathcal F_p\le2\sqrt3\,p^{3/2}.}
\end{equation}
In particular $\mathcal F_p\ll p^{3/2}$ unconditionally.
\end{theorem}

\begin{proof}
Partition $I_p$ into $J=\lceil M/H\rceil$ consecutive blocks of length at
most $H$, and let $m_j(v)$ count visits to $v$ in block $j$.  For each $v$,
Cauchy--Schwarz gives
\[
 m_p(v)^2=\left(\sum_{j=1}^Jm_j(v)\right)^2
 \le J\sum_{j=1}^Jm_j(v)^2.
\]
Within all blocks, the diagonal terms contribute $M$.  An unordered
off-diagonal collision in one block has a unique gap $2\le h\le H-1$ and
contributes twice to the squared multiplicity.  Therefore
\[
 \sum_{j,v}m_j(v)^2
 \le M+2\sum_{h=2}^{H-1}C_h
 \le M+6\sum_{q=1}^{H-2}q
 =M+3(H-1)(H-2).
\]
This proves~\eqref{eq:oe-block-bound}.  Notice that the factor $2$ for the
ordered pair is essential.

It remains to include the integer-rounding cost.  If $M\ge17$, take
$H=\lceil\sqrt{M/3}\rceil$.  Then
$\sqrt{M/3}\le H\le\sqrt{M/3}+1$ and
$\lceil M/H\rceil\le M/H+1$.  Expanding the right side of
\eqref{eq:oe-block-bound} and using these two inequalities gives
\begin{align*}
 \mathcal F_p
 &\le (M/H+1)(M+3H^2-9H+6)\\
 &\le2\sqrt3\,M^{3/2}-4M+8\sqrt{3M}+9
 \le2\sqrt3\,M^{3/2}.
\end{align*}
The last inequality holds for $M\ge17$ (the left side of
$4M-8\sqrt{3M}-9\ge0$ is already positive at $M=17$ and increasing).
Since $M<p$, this implies~\eqref{eq:oe-three-halves}.

For $5\le M\le16$, the trivial bound $\mathcal F_p\le M^2$ suffices:
$M^2/(M+2)^{3/2}$ is increasing, and at $M=16$ it is
$256/18^{3/2}<2\sqrt3$.  These are exactly the remaining primes $7\le p\le18$.
\end{proof}

\begin{corollary}[position relative to the pointwise bound]
\label{cor:oe-pointwise-position}
Theorem~\ref{thm:oe-energy-three-halves} is an energy theorem, not an
improvement of the pointwise exponent $2/3$.  More generally, an estimate
$\mathcal F_p\ll p^\theta$ improves $|Z_p|\ll p^{2/3}$ through energy alone
only when $\theta<4/3$.
\end{corollary}

\begin{proof}
Since $|Z_p|=m_p([0:1])$, one has $|Z_p|^2\le\mathcal F_p$.  Thus energy
gives exponent $\theta/2$, which is smaller than $2/3$ exactly when
$\theta<4/3$.
\end{proof}

% Ledger source: item 20 and §§90, 103.  This is deliberately fixed-H and
% prime-averaged; it is not the missing pointwise growing-family theorem.
\begin{proposition}[fixed-family Chebotarev mean]
\label{prop:oe-fixed-H-chebotarev}
Fix $H\ge1$.  Let $\nu_m$ be the number of distinct irreducible factors of
$U_m$ over $\mathbb Q$.  Apart from the finitely many primes at which one of
these polynomials has bad reduction,
\[
 \lim_{X\to\infty}\frac1{\pi(X)}
 \sum_{p\le X}\sum_{m=1}^H
 \#\{s\in\mathbb F_p:U_m(s)=0\}
 =\sum_{m=1}^H\nu_m.
\]
\end{proposition}

\begin{proof}
For one fixed squarefree polynomial, Frobenius fixes exactly as many roots
as the corresponding Galois-group element fixes in its action on the roots.
Chebotarev equidistributes the Frobenius conjugacy classes, and Burnside's
lemma says that the mean number of fixed roots is the number of Galois
orbits, namely the number of irreducible factors.  Apply this to the finitely
many polynomials $U_1,\ldots,U_H$ and sum.
\end{proof}

\begin{remark}[finite factor certificate]\label{rem:oe-factor-certificate}
The exact computations recorded in ledger §§90 and 109 certify, for
$m\le13$, that $\nu_m=2$ for odd $m$ and $\nu_m=1$ for even $m$, after the
shift $U_m(s)=N_{m+1}(s-1)$.  This is a finite machine certificate, not an
all-$m$ irreducibility theorem.
\end{remark}

\begin{remark}\label{rem:oe-chebotarev-limitation}
Proposition~\ref{prop:oe-fixed-H-chebotarev} averages over primes with $H$
fixed.  It neither controls $m$ growing like $\sqrt p$ at one prime nor
implies a pointwise bound for every prime.  That interchange of quantifiers
is the family-compatibility problem.
\end{remark}

% Ledger sources: §§103--104, 107--110; Q6506 §5; Q6511/Q6515 as indexed by
% DOSSIER_E1_EMPIRICAL.md.  Section §110 expressly corrects "equivalent" to
% "parallel sufficient/open targets"; that correction is enforced below.
\subsection{The sigma ladder and the E1 program}

The exact family form of energy is
\begin{equation}\label{eq:oe-energy-gap-identity}
 \mathcal F_p=M+2\sum_{h=1}^{M-1}C_h.
\end{equation}
Indeed, every unordered equal-color pair has one positive gap.

\begin{conjecture}[linear orbit energy, E1]\label{conj:oe-E1}
There is an absolute constant $C$ such that
\[
 \mathcal F_p\le Cp
\]
for every prime $p\ge7$.
\end{conjecture}

By~\eqref{eq:oe-energy-gap-identity}, E1 is equivalent to the linear family
bound $\sum_hC_h\ll p$.  This is the precise family-compatibility statement.

For $1\le\Delta\le M-1$, define
\[
 d_\Delta(r)=\#\{1\le h\le\Delta:r+h\le M,
                  \ \pi(r)=\pi(r+h)\},
 \qquad
 Q_p(\Delta)=\sum_r\binom{d_\Delta(r)}2.
\]

\begin{proposition}[the sigma exponent machine]\label{prop:oe-sigma-machine}
Fix $0\le\sigma\le2$.  Suppose that for every $\varepsilon>0$, uniformly
for primes $p$ and $1\le\Delta\le M-1$,
\begin{equation}\label{eq:oe-Qsigma}
 Q_p(\Delta)\ll_\varepsilon
 \Delta^{3-\sigma}p^\varepsilon.
\end{equation}
Then
\begin{equation}\label{eq:oe-sigma-output}
 \mathcal F_p\ll_\varepsilon
 p^{2-1/(3-\sigma)+\varepsilon}.
\end{equation}
Hence $\sigma=0,1,2$ correspond respectively to the exponents
$5/3,3/2,1$.  Theorem~\ref{thm:oe-energy-three-halves} occupies the
$\sigma=1$ rung, although its direct pair-gap proof does not assert the
stronger hypothesis~\eqref{eq:oe-Qsigma}.  At $\sigma=2$ the conclusion is
the exponent-level form $\mathcal F_p\ll_\varepsilon p^{1+\varepsilon}$;
the epsilon-free endpoint is Conjecture~\ref{conj:oe-E1}.
\end{proposition}

\begin{proof}
Put $S_p(\Delta)=\sum_rd_\Delta(r)$.  Since
$\sum_rd_\Delta(r)^2=S_p(\Delta)+2Q_p(\Delta)$, Cauchy--Schwarz gives
\[
 S_p(\Delta)^2\le M(S_p(\Delta)+2Q_p(\Delta)),
 \qquad
 S_p(\Delta)\le M+\sqrt{2MQ_p(\Delta)}.
\]
Partitioning $I_p$ into $\lceil M/\Delta\rceil$ blocks and repeating the
proof of Theorem~\ref{thm:oe-energy-three-halves} yields
\[
 \mathcal F_p\le
 \left\lceil\frac M\Delta\right\rceil
 \bigl(M+2S_p(\Delta)\bigr)
 \ll \frac{p^2}{\Delta}
 +p^{3/2+\varepsilon}\Delta^{(1-\sigma)/2}+p.
\]
Taking $\Delta$ to be the nearest admissible integer to
$p^{1/(3-\sigma)}$ balances the first two terms and proves
\eqref{eq:oe-sigma-output} (after renaming $\varepsilon$).
\end{proof}

We now state three concrete targets.  They are all motivated by the same
quenched, growing-gap family, but they are not mutually equivalent.

\begin{conjecture}[Ap\'ery triple return, ATR]
\label{conj:oe-ATR}
For $h\ne k$ put
\[
 J_p(h,k)=\#\{1\le r\le M-\max(h,k):
 \pi(r)=\pi(r+h)=\pi(r+k)\}.
\]
There is an absolute constant $C$ such that, for every prime $p\ge7$,
\begin{equation}\label{eq:oe-ATR}
 \sum_{\substack{1\le h,k\le M-1\\h\ne k}}J_p(h,k)\le Cp.
\end{equation}
This is an open conjecture, supported numerically but not proved.
\end{conjecture}

The implication to E1 is elementary.  If
$T_3=\sum_v(m_p(v))_3$, then the nonwrapping ordered convention gives
$3\sum_{h\ne k}J_p(h,k)=T_3$.  Moreover
$(m)_2\le(m)_3$ for $m\ge3$, while all fibres of size $2$ contribute at
most $M$ to $\sum_v(m_p(v))_2$.  Hence
\begin{equation}\label{eq:oe-ATR-implication}
 \mathcal F_p\le2M+T_3
 =2M+3\sum_{h\ne k}J_p(h,k),
\end{equation}
so Conjecture~\ref{conj:oe-ATR} implies E1.

Let $\kappa_h=1$ when $h$ is even and $\kappa_h=0$ when $h$ is odd.  The
pointwise reflection theorem supplies exactly this forced mirror collision.

\begin{conjecture}[growing-gap projective return variance, GPRV]
\label{conj:oe-GPRV}
There is an absolute constant $C$ such that, for every prime $p\ge7$,
\begin{equation}\label{eq:oe-GPRV}
 \sum_{h=1}^{p-3}
 \left|C_h-\kappa_h-
 \frac{p-2-h-\kappa_h}{p+1}\right|^2\le Cp.
\end{equation}
This is an open variance conjecture for the growing family of primitive
factors of the gap polynomials.
\end{conjecture}

By Cauchy--Schwarz,~\eqref{eq:oe-GPRV} makes the sum of the absolute
deviations $O(p)$.  The forced and uniform baseline terms also sum to
$O(p)$, so $\sum_hC_h=O(p)$; by~\eqref{eq:oe-energy-gap-identity}, GPRV
implies E1.

Finally put $v_n=(b_n,c_n)$ and
$\Delta_{r,h}=\det(v_r,v_{r+h})$.  For
\[
 \mathcal I_H=\{(r,h):1\le h\le H,\ 1\le r\le M-h\},
 \qquad e_p(x)=\exp(2\pi i x/p),
\]
define the genuinely two-parameter bilinear sum
\[
 B_t(H)=\sum_{(r,h)\in\mathcal I_H}e_p(t\Delta_{r,h}).
\]

\begin{conjecture}[two-dimensional square-root cancellation,
GAP--2D--SQRT]\label{conj:oe-2D-sqrt}
For every $\varepsilon>0$ there are $C_\varepsilon,p_\varepsilon$ such that,
for every prime $p\ge p_\varepsilon$, every $1\le H\le\sqrt p$, and every
$t\in\mathbb F_p^\times$,
\begin{equation}\label{eq:oe-2D-sqrt}
 |B_t(H)|\le C_\varepsilon(pH)^{1/2}p^\varepsilon.
\end{equation}
This is open; it asks for cancellation jointly in the base point and the
growing gap.
\end{conjecture}

Additive orthogonality gives the exact implication
\[
 \sum_{h\le H}C_h
 =\frac{\#\mathcal I_H}{p}
 +\frac1p\sum_{t\ne0}B_t(H)
 \ll_\varepsilon H+(pH)^{1/2}p^\varepsilon.
\]
Thus GAP--2D--SQRT breaks the elementary $H^2$ small-gap barrier (at
$H=\sqrt p$ it gives $p^{3/4+\varepsilon}$), but it does not by itself imply
E1.  This is why calling ATR, GPRV, and GAP--2D--SQRT ``equivalent'' would be
incorrect: the last is a parallel partial target, and the first two are
stronger sufficient criteria.

\begin{remark}[numerical status only]\label{rem:oe-empirical-status}
The computations collected in \texttt{DOSSIER\_E1\_EMPIRICAL.md} report
GPRV at the scale $\theta=0$, with the normalized variance near $1.08$ and
no visible drift over the tested range.  The dossier also reports agreement
to four decimal places for the principal energy, spectral, return, and
extreme-value laws of the mirror-random model.  These are numerical data,
not theorems, and none is used above.
\end{remark}

% Ledger sources: §108--109; the annealed comparison is the two-step return
% calculation recorded in theorem inventory item 4 and cited there to Q6456.
\subsection{The first splitting law and thin exceptional families}

Let
$R_2^{\rm alg}(p)=\#\{r\in\mathbb F_p:N_2(r)=0\}$ denote the full algebraic
root count (not the restricted-window count of
\eqref{eq:oe-Rrel}).

\begin{proposition}[the $h=2$ splitting law]\label{prop:oe-minus51}
For every prime $p\ge7$,
\begin{equation}\label{eq:oe-minus51}
 R_2^{\rm alg}(p)
 =1+2\,\mathbf 1_{\{(\frac{-51}{p})=1\}}.
\end{equation}
\end{proposition}

\begin{proof}
The dictionary gives
\[
 N_2(r)=P(r+1)
 =(2r+3)\bigl(17(r+1)^2+17(r+1)+5\bigr).
\]
The linear factor contributes one root.  The quadratic has discriminant
$17^2-4\cdot17\cdot5=-51$, so it contributes two roots when
$(-51/p)=1$ and none when $(-51/p)=-1$.  If $p\mid51$, the only case in
our range is $p=17$; the quadratic specializes to the nonzero constant $5$,
and the displayed formula again gives one root because the Legendre symbol
is zero.
\end{proof}

\begin{remark}[the same quadratic field in the annealed chain]
\label{rem:oe-minus51-annealed}
The independently proved annealed two-step return probability is
\[
 \frac{2+(\frac{-51}{p})}{p-1}.
\]
Thus the same field $\mathbb Q(\sqrt{-51})$ governs both the quenched
$h=2$ splitting law and the annealed two-step return.  This is the first
explicit warning that a proof of ATR, GPRV, or a related family estimate
must stratify fixed-field, thin exceptional subfamilies before applying a
generic random-family heuristic.
\end{remark}
